\section{\rm{VASP}软件的特色}
\subsection{\rm{VASP}的并行优势:~轮询调度算法}
\subsection{\rm{VASP}的精度保证:~双网格技术}

\section{\rm{VASP}的主要输出文件}
对用户而言，准备好输入文件并成功提交服务器，只是完成\textrm{VASP}计算的第一步。计算完成后，所有结果都包含在\textrm{VASP}产生的一系列输出文件中，如何从这些输出文件中快速、准确地提取体系能量、结构、电子态密度、能带、电荷密度等信息，是每一位\textrm{VASP}用户必须掌握的核心技能。

VASP的输出文件体系庞大而完整，主要包括：最主要的输出文件\textrm{OUTCAR}，其中包含计算过程的完整日志；结构文件\textrm{CONTCAR}，与\textrm{POSCAR}不同，\textrm{CONTCAR}记录了每次迭代后更新后的晶体结构；态密度文件\textrm{DOSCAR}，能量本征值文件\textrm{EIGENVAL}，电荷密度文件\textrm{CHGCAR}，体系本征态波函数文件\textrm{WAVECAR}，投影态波函数文件\textrm{PROCAR}，迭代计算过程中能量收敛纪录文件\textrm{OSZICAR}，局域势函数文件\textrm{LOCPOT}，电子局域函数文件\textrm{ELFCAR}，\textrm{Brillouin}区不可约$\vec k$点文件\textrm{IBZKPT}，离子轨迹文件\textrm{XDATCAR}等。此外，\textrm{VASP}还支持输出\textrm{XML}格式的输出文件\textrm{vasprun.xml}%和\textrm{HDF5}格式的输出文件\textrm{vaspout.h5}(需要有相关库函数支持)，
，为后续的数据处理和可视化提供了便捷的途径。

\subsection{主要输出文件:~\rm{OUTCAR}}
%# VASP输出文件完全解读：从OUTCAR到vaspout.h5

%## 引言
%本文将对VASP的各类输出文件进行全面、深入的解读，从文件格式、数据组织方式到实际应用中的读取与处理方法，力求为读者提供一份系统而详尽的参考指南。

%## 一、OUTCAR：VASP最主要的输出文件

%### 1.1 OUTCAR概述
\textrm{OUTCAR}是\textrm{VASP}计算中最重要的输出文件，它包含计算过程的主要信息，记录了从计算开始到结束的完整过程，包括汇总的输入参数、电子步自洽迭代{(Self-Consistent Field, SCF)}的每一步详细信息，以及离子步运动中原子在当前位置受力、应力张量、局域电荷与磁矩、介电性质等。\textrm{OUTCAR}是获取体系能量信息、\textrm{Fermi}能级以及原子受力信息的主要来源。\textrm{OUTCAR}中的信息详略，可以通过\textrm{INCAR}文件中的参数\texttt{NWRITE}标签进行控制:~\texttt{NWRITE}取值越大，\textrm{OUTCAR}输出的信息越详细，文件也越大。
%### 1.2 OUTCAR的文件结构
一般地，\textrm{OUTCAR}文件大致划分为几个主要部分，各部分之间用虚线隔开：
\begin{itemize}
	\item \textrm{VASP}的版本和计算资源信息；
	\item \textrm{POTCAR}中关于原子的基本信息；
	\item 根据计算模型\textrm{POSCAR}文件数据，确定最近邻原子间距离和晶体所属的对称性分析；
	\item 全部计算任务的控制参数，包括\textrm{INCAR}文件未设定的参数及其默认值；
	\item 计算任务过程的简要描述；
	\item 根据倒空间中$\vec k$及其权重等信息，计算内存分配等初始化计算信息；
	\item 电子步迭代信息，包括体系总能来源的各部分，\textrm{Fermi}能级和\textrm{Kohn-Sham}方程本征值等；
	\item 离子步计算的各方向应力张量；
	\item 离子步计算的每个原子的受力；
	\item 局域电荷与原子磁矩和总磁矩；
	\item 体系的介电性质。
\end{itemize}
%### 1.3 OUTCAR中的关键信息解析
%#### （一）输入参数汇总
\textrm{OUTCAR}的起始部分主要是计算软件和计算资源的基本信息。
{\fontsize{8.2pt}{6.2pt}\selectfont{
	\lstinputlisting[linerange=1-10]{Figures/OUTCAR} %为保险:~选用文件名绝对路径
}}
随后部分来自计算对象输入文件\textrm{INCAR}、\textrm{POTCAR}和\textrm{POSCAR}等的基本数据，主要是提供软件和用户校对和检查初始模型和参数是否一致(如\textrm{POSCAR}中的元素顺序与\textrm{POTCAR}中的元素顺序是否一致等)；此外，根据\textrm{POSCAR}中原子坐标和结构，检查晶体所属点群对称性，也是常规操作。在此不作过多说明。值得一提的是，如果用户有兴趣对\textrm{INCAR}所设置的参数及其内涵作系统了解，建议仔细看一下\textrm{OUOTCAR}中所列的完整的计算任务控制参数信息。
{\fontsize{8.2pt}{6.2pt}\selectfont{
	\lstinputlisting[linerange=215-366]{Figures/OUTCAR} %为保险:~选用文件名绝对路径
}}
这里根据参数在\textrm{VASP}计算过程中的作用与分类，详细罗列各参数的取值与作用，这部分内容是用户熟悉\textrm{INCAR}参数的最有效的途径。

在确定所有参数后，\textrm{VASP}软件将提示用户本次计算主要完成的任务。
{\fontsize{8.2pt}{6.2pt}\selectfont{
	\lstinputlisting[linerange=367-388]{Figures/OUTCAR} %为保险:~选用文件名绝对路径
}}
依次说明计算类型，如只有电子步计算的静态计算\textrm{(static calculation)}，或包含离子运动的结构弛豫\textrm{(relaxation of ions)}或分子动力学计算；电子步迭代过程中电荷密度与势函数是否变化；是否考虑自旋极化；矩阵对角化方案；电荷密度的混合方案；应力计算方案；电荷密度积分(\textrm{Fermi}面计算)是否引入\textrm{Gaussian}展宽等等。这些信息是对本次计算的总体概括。

此后，\textrm{VASP}根据计算要求进入计算条件的初始化阶段，主要是分配计算内存等，这部分内容主要是计算资源分配，用户可以忽略。然后进入\textrm{VASP}计算的主体，即离子步和电子步计算迭代流程。
{\fontsize{8.2pt}{6.2pt}\selectfont{
	\lstinputlisting[linerange=1011-1058]{Figures/OUTCAR} %为保险:~选用文件名绝对路径
}}
\textrm{OUTCAR}以离子步-电子步迭代方式记录完整的迭代计算过程，每个完成的迭代计算起始以\textrm{Iteration~A(B)}为标记，其中\textrm{A}表示离子步序号，括号中的\textrm{B}表示该离子步对应的电子步序号。以下汇集该离子步计算的各部分用时统计和电荷密度混合收敛判断。电子步迭代最重要的输出是相关能量数据，体系计算自由能\textrm{Free energy}的各部分分解，是按照赝势方法下体系总能计算公式分配的，\textcolor{red}{\textbf 对这部分内容由兴趣的读者，可以参阅本书关于体系总能的计算公式。}

\textrm{VASP}计算的体系总能，分别以
\begin{itemize}
	\item 自由能\textrm{free energy TOTEN}
	\item 扣除熵贡献的能量\textrm{energy without entropy}
\end{itemize}
形式给出。这两者的差别，根源于对最高占据态的\textrm{Fermi}面积分函数的近似，原则上\textrm{DFT}计算的计算的是\textrm{T=0K}时的电子态，但这种情况下，总能量计算中对电子最高占据态的积分将会是$\delta$函数，所以一般采用带展宽的近似函数(如\textrm{Fermi}分布函数或\textrm{Gaussian}函数)近似替代，这也是为什么\textrm{INCAR}文件中的控制参数有\texttt{SIGMA}的原因。当最高占据态积分采用四面体积分方案(\texttt{ISMERA}=-5)时，没有展宽参数，此时体系自由能即为总能；而当$\mathtt{ISMERA}\neq-5$，则``扣除熵贡献的能量''对应体系总能，计算结合能等相关能量时，需要采用该值。

每一个离子步结束，体系会输出\textrm{Fermi}能和$\vec k$点的能量本征值和轨道占据数:
{\fontsize{8.2pt}{6.2pt}\selectfont{
	\lstinputlisting[linerange=1153-1172]{Figures/OUTCAR} %为保险:~选用文件名绝对路径
}}
在每个离子步内，最后一个电子步结束后，\textrm{VASP}会提示``\textrm{aborting loop because EDIFF is reached}''，然后提供电荷密度、原子应力和受力等各类物理量的计算时间：
{\fontsize{8.2pt}{6.2pt}\selectfont{
	\lstinputlisting[linerange=1197-1208]{Figures/OUTCAR} %为保险:~选用文件名绝对路径
}}
然后\textrm{VASP}将输出的是原子在各个方向上的应力\textrm{(Stress)}和应力张量\textrm{(Stress Tensor)}，原子受力\textrm{(Force)}以及更新后的晶胞参数和原子位置(晶格和原子位置许可的变化方式，由参数\textrm{ISIF}控制并确定)：
{\fontsize{8.2pt}{6.2pt}\selectfont{
	\lstinputlisting[linerange=1209-1266]{Figures/OUTCAR} %为保险:~选用文件名绝对路径
}}
如果离子弛豫最后达到收敛，则\textrm{VASP}输出提示``\textrm{reached required accuracy - stopping structural energy minimisation}''，并提示波函数等的输出，以及内存空间的释放情况；如果是分子动力学计算，则在达到离子变化步数上限后停止，后续输出则类似。此外如有其它物理性质(如磁学、光学性质等)也将在此输出。
{\fontsize{8.2pt}{6.2pt}\selectfont{
	\lstinputlisting[linerange=2951-2968]{Figures/OUTCAR} %为保险:~选用文件名绝对路径
}}
全部计算结束，\textrm{VASP}在结尾总结全部计算的耗时、内存使用情况等汇总信息:
{\fontsize{8.2pt}{6.2pt}\selectfont{
	\lstinputlisting[linerange=2969-2983]{Figures/OUTCAR} %为保险:~选用文件名绝对路径
}}
%会\textrm{INCAR}中的所有参数逐行打印出来，方便用户确认计算设置是否正确。这部分还包括从\textrm{POTCAR}中读取的赝势信息、从\textrm{POSCAR}中读取的结构信息以及从\textrm{KPOINTS}中读取的k点设置
%#### （二）能量信息
%\textrm{OUTCAR}中包含两类重要的能量值：
%
%- **free energy TOTEN**：自由能
%- **energy without entropy**：不含熵项的能量
%
%这两者的选择取决于INCAR中的**ISMEAR**设置：
%- 当`ISMEAR = -5`（四面体方法）时，两者相等，可通过`grep 'TOTEN' OUTCAR`获取能量数据。
%- 当`ISMEAR ≠ -5`时，两者不相等，计算结合能时应选取**energy without entropy**对应的数值，可通过`grep 'entropy=' OUTCAR`获取。
%
%此外，费米能级可通过`grep 'Fermi' OUTCAR`获取。
%
%%#### （三）电子自洽迭代信息
%在电子SCF循环中，OUTCAR会输出每一步的迭代信息：
%- 当前步数
%- 总能量变化
%- 电荷密度变化
%- 收敛状态
%
%当电子步收敛时，OUTCAR会显示`reached required accuracy`或类似信息。
%
%%#### （四）离子步信息
%对于结构优化（`IBRION=1或2`）或分子动力学（`IBRION=0`）计算，OUTCAR会输出每个离子步的信息：
%- 离子步编号
%- 原子位置
%- 原子受力（forces）
%- 晶格参数变化（若`ISIF`允许）
%
%原子受力信息可通过`grep 'TOTAL-FORCE' OUTCAR`定位。收敛标准（`EDIFFG`）是否满足也会在OUTCAR中明确标示。
%
%%#### （五）应力张量
%当`ISIF > 0`时，OUTCAR会输出应力张量（stress tensor），单位为kBar。这对于研究体系在压力下的行为或进行晶胞优化至关重要。
%
%### 1.4 如何高效读取OUTCAR
%
\textrm{OUTCAR}的文件内容通常较为详细，根据计算对象原子数的不同，文件大小从\textrm{MB}到\textrm{GB}量级不等。对于计算规模较大的体系，通览全部文件必要性不大。可以借助\textrm{Linux}的\textrm{Bash}脚本(\textrm{Script})命令功能或\textrm{VASPKIT}等工具软件，快速提取计算过程中的关键物理量或相关信息。
%
%| 目的 | 命令 |
%|------|------|
%\begin{itemize}
%		\item 获取能量信息:~ \texttt{grep 'TOTEN' OUTCAR}或\textt{grep 'entropy=' OUTCAR}
%		\item 获取\textrm{Fermi}能级信息:~\texttt{grep 'E-fermi' OUTCAR}
%		\item 获取原子受力信息:~ `grep 'TOTAL-FORCE' OUTCAR` |
%| 获取应力张量 | `grep 'stress' OUTCAR` |
%| 检查是否收敛 | `grep 'reached required' OUTCAR` |
%| 获取磁矩 | `grep 'magnetization' OUTCAR` |
%\end{itemize}
%
%### 1.5 OUTCAR的常见问题
除了\textrm{OUTCAR}文件，\textrm{VASP}还会以\textrm{XML}格式形式将结果输出到\textrm{vasprun.xml}文件中，它有着比\textrm{OUTCAR}更结构化的呈现形式。\textrm{vasprun.xml}的优势主要在于:
\begin{itemize}
	\item 便于解析:~\textrm{XML}格式便于程序化读取;
	\item 与可视化工具\textrm{p4vasp}兼容:~\textrm{p4vasp}可以直接读取\textrm{vasprun.xml}数据将结果可视化。
\end{itemize}

%1. **文件过大**：对于大型计算，OUTCAR可能达到数GB。可通过设置`NWRITE = 1`（默认）或`NWRITE = 0`减少输出量。
%2. **能量选择错误**：如前所述，不同ISMEAR设置下应选择不同的能量值。
%3. **收敛判断**：OUTCAR中可能出现“电子步未收敛”的警告，此时需要检查INCAR中的`NELM`、`ALGO`等参数。

\subsection{其余输出文件}
\textrm{VASP}的主要计算结果都体现在\textrm{OUTCAR}文件中，但有些重要的结果数据，如结构弛豫后的晶体和原子数据、用于绘制相关物理量图表的数据，\textrm{VASP}会将必要的数据单独存储到特定文件中。这些输出文件中的数据功能相对具体，这里简要地介绍几个文件。
%## 二、CONTCAR：更新后的结构文件
%
%### 2.1 CONTCAR概述
%

\textrm{CONTCAR}是\textrm{VASP}在每个离子步结束后记录当前计算对象的晶体结构和原子位置的文件，其文件格式与\textrm{POSCAR}完全相同，因此如果离子步计算中断，可以直接复制为\textrm{CONTCAR}的数据到\textrm{POSCAR}，重启\textrm{VASP}进行后续计算。显然，整个计算任务结束时，\textrm{CONTCAR}保存的就是体系最终能达到的晶体结构和离子能达到的位置。
\begin{itemize}
	\item 对于静态计算(\texttt{IBRION=-1}):~\textrm{CONTCAR}与初始\textrm{POSCAR}内容相同。
	\item 对于结构优化计算(\texttt{IBRION>1}):~\textrm{CONTCAR}包含结构优化后的最终结构。
	\item 对于\textrm{MD}模拟(\texttt{IBRION=0}):~\textrm{CONTCAR}除了与\textrm{POSCAR}中的信息外，还包含原子的初始速度\textrm{(velocites)}和预测-校正坐标\textrm{(predictor-corrector coordinates)}等数据。这些数据可用于后续的\textrm{MD}计算，需要指出的是，如果用户从\textrm{CONTCAR}数据开始继续\textrm{MD}计算，但更换系综，比如从\textrm{NVT}系综切换到\textrm{NPT}系综，必须移除预测-校正坐标数据，否则会出现计算不兼容的现象。
\end{itemize}
%
%CONTCAR的核心作用是**保存计算结束后的原子位置**——
%
%### 2.2 CONTCAR的文件格式
%
%CONTCAR的格式与POSCAR完全一致，包含以下内容：
%- **第1行**：注释行
%- **第2行**：全局缩放因子
%- **第3-5行**：晶格向量
%- **第6行**（VASP 5.x及以上）：原子种类符号
%- **第7行**：各元素原子数
%- **第8行**：坐标类型（Direct或Cartesian）
%- **第9行及之后**：原子坐标
%
%### 2.3 分子动力学中的CONTCAR
%
%对于分子动力学模拟（`IBRION=0`），CONTCAR包含的内容更加丰富：
%- **第一块**：晶格参数和原子坐标
%- **第二块**：原子的初始速度（velocities）
%- **第三块**：预测-校正坐标（predictor-corrector coordinates）
%
%这些额外信息使得CONTCAR可以直接用于继续MD计算。
%
%**重要提醒**：如果要从CONTCAR继续MD计算但更换系综（如从NVT切换到NPT），必须移除第三块的预测-校正坐标，否则计算将不兼容。
%
%### 2.4 CONTCAR的使用场景
%
%1. **结构优化后续计算**：将CONTCAR复制为POSCAR，进行更高精度的静态计算或性质计算。
%2. **继续未完成的结构优化**：若优化因时间限制中断，可直接用CONTCAR继续。
%3. **分子动力学续算**：CONTCAR包含了继续MD所需的所有信息。
%
%### 2.5 CONTCAR的格式注意事项
%
%- CONTCAR中每行开头的缩进使用的是**空格**而非制表符（Tab）。
%- VASP 5.x版本在晶格向量之后增加了原子种类符号行，这使得CONTCAR的向后兼容性更好。

\textrm{DFT}计算过程中，用户主要关心的是体系能量的变化趋势。\textrm{VASP}计算过程中，频繁地检索逐渐增长的\textrm{OUTCAR}文件来监控中间计算过程，显然不是最便捷的方法，而\textrm{OSZICAR}文件\footnote{文件名中的\textrm{``OSZI''}源自德语\textrm{``oszillieren''}，这个词是``振荡''的意思。}就是计算过程中简要的能量汇总文件(如果用户直接运行\textrm{vasp}执行文件，在屏幕上显示的主要也是\textrm{OSZICAR}文件的内容)。典型的\textrm{OSZICAR}文件的内容如下所示:
{\fontsize{8.2pt}{6.2pt}\selectfont{
	\lstinputlisting[linerange=1-24]{Figures/OSZICAR} %为保险:~选用文件名绝对路径
}}
按列划分，\textrm{OSZICAR}的内容依次为:
\begin{itemize}
	\item 电子步迭代对角化方法；
	\item \textrm{N}:~电子步迭代步数编号；
	\item \textrm{E}:~当前电子步计算得到的体系自由能；
	\item \textrm{dE}:~体系当前与前一步的自由能的变化，用于判断收敛趋势；
	\item \textrm{d eps}:~体系的带结构能量(本征值之和$\sum\limits_{\vec k,i}^\mathrm{occ}\varepsilon_{\vec k,i}$)变化；
	\item \textrm{ncg}:~矩阵迭代对角化过程中\textrm{Hamiltonian}矩阵作用于波函数的评估次数；
	\item \textrm{rms (c)}:~\textrm{rms}表示残矢$\mathbf{R}=(\mathbf{H}-\varepsilon\mathbf{S})|\phi\rangle$的模量，\textrm{rms (c)}则是电荷密度迭代的残差的均方根。
\end{itemize}
最后给出的是关于电子步收敛的体系总能的参考信息，包括第一项是总自由能\textrm{F}，第二项$\mathrm{E}_0$是\textrm{OUTCAR}中\textrm{energy(sigma->0)}对应的能量，第三项\textrm{d E}是当前步和前一步总能的差(对于静态计算，此处\textrm{d E}是\textrm{OUTCAR}中的熵\textrm{(entropy T*S EENTROP)})。%与参数\textrm{SIGMA}的乘积)
\textrm{MD}计算中，最后这部分各项略不同:
{\fontsize{8.2pt}{6.2pt}\selectfont{
	\lstinputlisting[linerange=16-16]{Figures/OSZICAR_MD} %为保险:~选用文件名绝对路径
}}
此处第一项\textrm{T}为当前离子步的温度；\mathrm{E}是总的自由能(包括离子动能和\textrm{Nos\'e}恒温方案的能量的贡献)；\textrm{F}与$\mathrm{E}_0$与前面一样；\textrm{EK}是体系的动能；\textrm{SP}是来自\textrm{Nos\'e}恒温方案的势能部分的贡献；\textrm{SK}则是来自\textrm{Nos\'e}恒温方案的动能部分贡献。

利用\textrm{OSZICAR}文件，用户可以在计算运行过程中，对体系能量变化有总体把握，如通过命令\texttt{tail -f OSZICAR}实时查看当前电子步能量收敛情况；通过命令\textrm{grep -in ' F= ' OSZICAR}实时查看离子步能量变化状况。
%2. **快速检查**：
计算结束后，\textrm{OSZICAR}文件可以帮助用户快速确认体系是否达到收敛及了解最终的能量状况。
%3. **机器学习力场**：当使用机器学习力场时，OSZICAR同样会被写入。

%### 8.1 OSZICAR概述
%**OSZICAR**文件记录了**电子自洽迭代**（SCF）和**离子步**过程中的能量变化信息。文件名中的"OSZI"源自德语"oszillieren"（振荡），"CAR"表示卡片/记录。OSZICAR提供了计算过程的**简洁摘要**，非常适合用于**实时监控**计算进度。
%
%### 8.2 OSZICAR的内容
%
%OSZICAR包含以下信息：
%- 选择的SCF算法
%- 总能量、电荷密度和自旋密度的收敛情况
%- 自由能（free energy）
%- 体系的磁矩
%
%### 8.3 OSZICAR的格式示例
%
%典型的OSZICAR内容如下：
%
%```
%N   E         dE        d eps    ncg   rms    rms(c)
%CG : 1  -.13238703E+04  -.132E+04  -.934E+02  56  .28E+02
%CG : 2  -.13391360E+04  -.152E+02  -.982E+01  82  .54E+01
%CG : 3  -.13397892E+04  -.653E+00  -.553E+00  72  .13E+01  .14E+00
%...
%1 F= -.13401522E+04 E0= -.13397340E+04 d E = -.13402E+04
%```
%
%各列含义如下：
%- **N**：电子步编号
%- **E**：当前自由能
%- **dE**：自由能变化
%- **d eps**：能带结构能量变化
%- **ncg**：哈密顿量作用于波函数的评估次数
%- **rms**：残差的均方根
%
%### 8.4 OSZICAR的实用价值
%
%1. **实时监控**：在计算运行过程中，通过`tail -f OSZICAR`可以实时查看能量收敛情况。
%2. **快速检查**：计算完成后，通过OSZICAR可以快速确认是否收敛及最终能量。
%3. **机器学习力场**：当使用机器学习力场时，OSZICAR同样会被写入。
%
%

%
%
%## 三、DOSCAR：态密度文件
%
%### 3.1 DOSCAR概述
%
态密度\textrm{(Density of States, DOS)}是呈现电子结构信息最有效的手段之一。\textrm{DOSCAR}是提供绘制总态密度\textrm{(TDOS)}和分波态密度\textrm{(PDOS)}图的数据来源。\textrm{DOSCAR}文件中的数据包含了体系的态密度和积分态密度\textrm{(Integrated DOS)}，相应的单位分别为``状态数\textrm{/eV}''和``状态数''。
%
%对于动态模拟（分子动力学）和结构弛豫计算，DOSCAR中写入的是**平均后的DOS和积分DOS**。
%
%### 3.2 DOSCAR的文件结构
%
\textrm{DOSCAR}文件的内容大致分为三个部分:
\begin{itemize}
	\item 文件信息描述 
{\fontsize{8.2pt}{6.2pt}\selectfont{
	\lstinputlisting[linerange=1-6]{Figures/DOSCAR} %为保险:~选用文件名绝对路径
}}
用户可以根据经验判断这些数据的含义，分别是离子数，是否包括\textrm{PDOS}(0或1)\\
晶体体积\textrm{$\AA^3$}\\
体系名称\\
能量范围(\textrm{EMAX},\textrm{EMIN})，\textrm{NEDOS}(\textrm{DOS}的数据点数，默认值为301)，\textrm{E-fermi}，固定值1.000\\
按照传统习惯，绘制\textrm{DOS}图时，\textrm{E-fermi}能置为能量零点。
%
%#### （一）文件头（Header）
%DOSCAR的前几行构成文件头，包含以下信息：
%- 离子总数（包括空球）
%- 是否包含分波DOS（0或1）
%- 晶胞体积（ų）
%- 晶格基矢长度（a, b, c，单位m）
%- POTIM（时间步长，单位s）
%- 初始温度（TEBEG）
%- 体系名称（SYSTEM）
%- 能量范围（EMAX, EMIN）
%- NEDOS（DOS数据点数）
%- 费米能级（Efermi）
%- 固定值1.0000
%
%#### （二）总DOS和积分DOS
\item \textrm{TDOS}和\textrm{Integrated DOS}数据:\\
%紧接着文件头的是**NEDOS行**数据，
	每行包含三列，依次为能量-\textrm{DOS}-\textrm{Integrated DOS}的数值。
%1. **能量**（Energy）
%2. **DOS**（态密度）
%3. **积分DOS**（Integrated DOS）
%
	对于自旋极化的计算\texttt{ISPIN=2}，\textrm{DOS}部分会分列出\textrm{spin-up}和\textrm{spin-dn}两部分。
%
%#### （三）分波DOS（可选）
\item 分波DOS(当$\mathtt{LORBIT}\geqslant10$)
	\textrm{DOSCAR}还会列出属于每个原子的\textrm{PDOS}。这部分数据以19列的形式依次列出，包含了每个原子各轨道(\textit{s}，\textit{p}，\textit{d}}等)的投影态密度。
%
%### 3.3 影响DOSCAR输出的INCAR参数
\end{itemize}
%
%以下INCAR参数直接影响DOSCAR的内容：
%
%| 参数 | 作用 |
%|------|------|
%| **NEDOS** | DOS数据点数，默认301 |
%| **EMIN** | DOS能量范围下限 |
%| **EMAX** | DOS能量范围上限 |
%| **NBLOCK** | 控制DOS平均的块大小 |
%| **KBLOCK** | 控制DOS平均的k点块大小 |
%| **LORBIT** | 控制是否输出分波DOS以及输出格式 |
%
%### 3.4 如何处理DOSCAR
%
\textrm{DOSCAR}的原始数据简单明了，通常可使用以下工具进行可视化：
\begin{itemize}
	\item \textrm{p4vasp}:~\textrm{VASP}官方推荐的图形化工具，可直接读取\textrm{DOSCAR}数据，并绘制\textrm{DOS/PDOS}图。
	\item \textrm{VASPKIT}:~可提取\textrm{DOSCAR}中的\textrm{DOS}和各原子\textrm{PDOS}数据，用以绘图。
%- **split_dos**：第三方脚本，可将DOSCAR分割为总态密度（DOS0）和各个原子的分波态密度（DOS1, DOS2, …）。
\end{itemize}
%
%使用split_dos时，需要保证当前目录下有对应的OUTCAR和POSCAR文件。
%
%### 3.5 DOSCAR的注意事项
%
%1. **费米能级**：DOSCAR中的能量是相对于费米能级的。绘制DOS图时通常将费米能级设为能量零点。
%2. **积分态密度**：积分态密度在费米能级处的值应等于体系的总电子数（对于绝缘体/半导体）。
%3. **四面体方法**：当使用`ISMEAR = -4`或`-5`（四面体方法）时，积分态密度的计算采用特殊公式。
%
%
%## 四、EIGENVAL：本征值文件
%
%### 4.1 EIGENVAL概述
%

除了态密度外，每个$\vec k$的能量本征值也是电子结构最重要的信息之一。\textrm{VASP}计算中所有$\vec k$点的求解\textrm{Kohn-Sham}方程得到的电子本征态能量(可用于绘制能带图的数据)都存于文件\textrm{EIGENVAL}中，因此计算结束，它保存的是最后一个离子步的最终电子能带数据。有必要指出的是，由于\textrm{VASP}程序的设定，在非静态计算时($\mathtt{IBRION}\neq-1$)，电子步求解的本征值的方法和精度都粗糙一些，因此\textrm{EIGENVAL}保存的电子结构数据，不建议直接用作电子结构的解析；如果需要当前结构(经历结构弛豫或\textrm{MD}计算)的电子结构数据，一个稳妥的做法是，先将\textrm{CONTCAR}复制为\textrm{POSCAR}，然后额外完成一次静态续算(\texttt{ISTART}=1，\texttt{NSW}=0，\texttt{ICHARG}=1，\texttt{IBRION}=-1)即可。
%
%### 4.2 不同计算类型中的EIGENVAL
%
%EIGENVAL的内容因计算类型而异：
%
%- **静态计算和结构弛豫**（`IBRION = -1, 1, 2`）：EIGENVAL包含最后一个离子步的KS方程解。
%
%- **分子动力学模拟**（`IBRION = 0`）：EIGENVAL包含的是**预测的下一步**的本征值，与CONTCAR文件兼容。**重要提醒**：如果要将MD中的EIGENVAL用于其他计算，需要先将CONTCAR复制为POSCAR，然后做一个静态续算（`ISTART=1, NSW=0, ICHARG=1`）。
%
%### 4.3 EIGENVAL的文件格式
%
%EIGENVAL的格式大致如下：
%- 文件头：包含体系信息、k点总数、能带总数等
%- 对每个k点：
%  - k点坐标和权重
%  - 该k点处所有能带的本征值（每个能带一行）
%
%### 4.4 EIGENVAL的用途
%
%EIGENVAL主要用于：
%1. **绘制能带结构图**：结合KPOINTS文件中的高对称路径，提取特定路径上的本征值绘制能带。
%2. **计算有效质量**：通过能带在带边附近的色散关系提取有效质量。
%3. **分析能隙**：确定价带顶和导带底的位置及能隙大小。
%
%
%## 五、CHGCAR：电荷密度文件
体系的电荷密度是\textrm{DFT}计算的核心，\textrm{VASP}计算中\textrm{CHGCAR}文件保存的就是这类数据。
\textrm{CHGCAR}文件主要由以下数据构成：
\begin{itemize}
	\item 计算对象描述:~参照\textrm{POSCAR}类似的形式，描述晶格参数和原子坐标：
{\fontsize{8.2pt}{6.2pt}\selectfont{
	\lstinputlisting[linerange=1-11]{Figures/CHGCAR} %为保险:~选用文件名绝对路径
}}
\item \textrm{FFT}网格划分与电荷密度:~\textrm{NGXF}`, `NGYF`, `NGZF`三个整数。
{\fontsize{8.2pt}{6.2pt}\selectfont{
	\lstinputlisting[linerange=12-15]{Figures/CHGCAR} %为保险:~选用文件名绝对路径
}}
	电荷密度数据以多行实数形式输出，共计$\mathrm{NGXF}\times\mathrm{NGYF}\times\mathrm{NGZF}$个数据点。
\item \textrm{USPP/PAW}方法中补偿电荷\textrm{(Augmentation occupancies)}，它只在每个原子附近的实空间里有贡献。
{\fontsize{8.2pt}{6.2pt}\selectfont{
	\lstinputlisting[linerange=23-38]{Figures/CHGCAR} %为保险:~选用文件名绝对路径
}}
\end{itemize}
\textrm{CHGCAR}保存的是自洽迭代中每一个电子步的电荷密度，因此其最重要的功用，就是在\textrm{VASP}重启计算时，可以通过设置\texttt{ICHARG}=1令程序直接从已有的\textrm{CHGCAR}文件中读取电荷密度进行续算。而避免从头开始重新计算，节约计算时间和成本。%，官方建议在反复重启计算（如微调k点网格）时，从CHGCAR开始。
而对于原子数较多的体系，\textrm{CHGCAR}文件可能过大，用户也可以通过设置参数\texttt{LCHARG}=\textrm{.FALSE.}避免数据写出到\textrm{CHGCAR}文件中。
%### 5.1 CHGCAR概述
%
%**CHGCAR**文件存储了体系的**电荷密度**以及**PAW的一中心占据数**（one-center occupancies）。它是VASP默认输出的文件之一（可通过设置`LCHARG = .FALSE.`避免输出，或通过`LH5`重定向到vaspwave.h5）。
%
%CHGCAR最重要的用途是**用于重启VASP计算**——通过设置`ICHARG = 1`从已有的CHGCAR读取电荷密度。官方建议在反复重启计算（如微调k点网格）时，从CHGCAR开始。
%
%### 5.2 CHGCAR的文件结构
%
%CHGCAR由以下几个数据块组成：
%
%1. **结构块**：采用与POSCAR完全相同的格式，包含晶格参数和原子坐标。
%2. **FFT网格维度**：`NGXF`, `NGYF`, `NGZF`三个整数。
%3. **电荷密度数据**：`电荷 × FFT网格体积`的值，以多行实数形式输出，直到写满`NGXF × NGYF × NGZF`个数据点。
%4. **增强占据数**（Augmentation occupancies）。
%
%### 5.3 CHGCAR中的电荷密度值
%
%CHGCAR中存储的电荷密度值是**电荷密度ρ(r)乘以网格体积**后的结果。因此，对CHGCAR中的所有数据点求和即可得到体系的总电子数。
%
%### 5.4 CHGCAR与CHG的区别
%
\textrm{VASP}还有名为\textrm{CHG}的输出文件，初学者不太清楚它与\textrm{CHGCAR}的关系。简言之，\textrm{CHG}主要提供的电荷密度可视化的绘图数据，其文件起始部分与\textrm{CHGCAR}相同，都是关于计算对象的结构和原子信息，但随后的电荷密度信息，是从第一个离子步起，每满十个离子步才往\textrm{CHG}文件一次。因此这些数据往往与\textrm{CHGCAR}的数据点没有对应关系。
%- **CHGCAR**：包含PAW一中心占据数，用于重启计算。
%- **CHG**：不包含PAW一中心占据数，主要用于**可视化**和后处理。
%
%### 5.5 CHGCAR的实用技巧
%
%1. **电荷密度差分**：通过比较不同体系的CHGCAR，可以计算电荷密度差分（charge density difference），用于分析成键和电荷转移。
%2. **平面平均**：对于表面和界面体系，可对CHGCAR做平面平均（planar average）以获得沿某一方向的电荷分布。
%3. **可视化**：CHGCAR可通过VESTA等软件直接可视化。
%

%## 六、WAVECAR：波函数文件
根据\textrm{DFT}理论，\textrm{Kohn-Sham}方程的波函数是中间变量。一般地，第一原理计算软件中波函数\footnote{因为选定了基函数，所以波函数文件中存储的，实际上是构成波函数的基组系数。}都是以非格式(二进制)输出形式写入文件的，用户必须借助工具才能读取数据。\textrm{VASP}也不例外，\textrm{WAVECAR}是一个二进制文件，也是\textrm{VASP}最大的输出文件之一。与\textrm{CHGCAR}类似，保存\textrm{WAVECAR}文件的主要作用之一，就是节省重启/续算时的计算时间。因为\textrm{WAVECAR}保存的是程序结束前最后完成的电子步的波函数，所以在续算任务启动时，用户如使用已有的\textrm{WAVECAR}数据(参数\texttt{ISTART}=1)，则有可能有效加速电子步\textrm{SCF}的收敛速率。

%### 6.1 WAVECAR概述
%
%**WAVECAR**是一个**二进制文件**，包含了体系的**波函数系数**及相关信息。它是VASP中体积最大的输出文件之一，通常可以用于从之前的计算继续。
%
%### 6.2 WAVECAR的内容
%
%WAVECAR文件中包含以下数据：
%- **NBAND**：能带总数
%- **ENCUTI**：初始截断能
%- **AX**：定义超胞的初始基矢
%- **CELEN**：初始本征值
%- **FERWE**：初始费米权重
%- **CPTWFP**：初始波函数
%
%### 6.3 WAVECAR的技术细节
%
%WAVECAR是一个Fortran二进制文件，其格式与编译时的精度设置有关：
%- `RTAG = 45200`：使用complex*8（单精度）格式
%- `RTAG = 45210`：使用complex*16（双精度）格式
%
%WAVECAR的总记录数为：`2 + NSPIN × NKPTS × (1 + NBANDS)`。
%
%### 6.4 不同计算类型中的WAVECAR
%
%- **静态计算和结构弛豫**（`IBRION = -1, 1, 2`）：WAVECAR包含最后一个离子步的KS方程解。
%- **分子动力学模拟**（`IBRION = 0`）：WAVECAR包含**预测的波函数**，与CONTCAR兼容。
%
%### 6.5 WAVECAR的控制
%
%可以通过在INCAR中设置`LWAVE = .FALSE.`来禁止输出WAVECAR文件。对于大型体系，关闭WAVECAR输出可以节省大量存储空间。
%
%### 6.6 WAVECAR的使用建议
%
%1. **续算加速**：使用已有的WAVECAR（`ISTART = 1`）可以显著加速SCF收敛。
%2. **MD续算**：WAVECAR、CHGCAR和CONTCAR可以一致地用于分子动力学续算。
%3. **注意一致性**：如果要从MD的WAVECAR进行其他计算，需要先做静态续算。
%

%## 七、PROCAR：投影波函数文件
%

在对\textrm{VASP}计算结果解析中，会有分析体系本征态(能带)波函数中各组成元素的\textit{s}、\textit{p}、\textit{d}\、\textit{f}等轨道的贡献多寡的需求，而\textrm{PROCAR}文件中听过的就是这些数据信息(要求控制参数$\mathtt{LORBIT}\geqslant11$)\footnote{在并行版本中，如果$\mathtt{NPAR}\neq1$`，则\textit{s}，\textit{p}，\textit{d}等投影的波函数数据不会被计算。}。
{\fontsize{8.2pt}{6.2pt}\selectfont{
	\lstinputlisting[linerange=1-23]{Figures/PROCAR} %为保险:~选用文件名绝对路径
}}
对与每个$\vec k$点上的每个本征态(能带)波函数，\textrm{PROCAR}输出一个数据块信息，包含了每个离子的各轨道投影值的贡献，据此，用户可以判断特定能带的电子主要来自哪些原子的哪些轨道，通过投影值的大小，也可以初步估计本征态(能带)的局域化程度。
%### 7.1 PROCAR概述
%
%**PROCAR**文件包含了每条能带的**spd轨道投影**和**位点投影**的波函数特征。它是分析电子态**轨道成分**和**局域化性质**的关键文件。
%
%### 7.2 PROCAR的生成条件
%
%PROCAR的生成需要满足以下条件：
%- 在INCAR中设置**LORBIT ≥ 10**或设置所有原子的**RWIGS**值。
%- 对于`LORBIT < 10`，必须在INCAR中显式指定`RWIGS`。
%
%**重要提醒**：
%
%### 7.3 PROCAR的文件格式
%
%PROCAR的文件格式如下：
%
%```
%# of k-points: 5
%# of bands: 26
%# of ions: 3
%k-point 1 : 0.00000000 0.00000000 0.00000000 weight = 0.06250000
%band 1 # energy -17.37867948 # occ. 1.00000000
%ion    s     py    pz    px    dxy   dyz   dz2   dxz   x2-y2   tot
%1   0.144 0.000 0.000 0.000 0.000 0.000 0.000 0.000 0.000 0.145
%2   0.291 0.000 0.006 0.000 0.000 0.000 0.000 0.000 0.000 0.297
%...
%```
%
%
%### 7.4 PROCAR的用途
%
%1. **轨道成分分析**：判断特定能带的电子主要来自哪些原子和哪些轨道。
%2. **能带标定**：在能带图上用颜色或点的大小标示轨道成分（“fatband”图）。
%3. **局域化分析**：通过投影值的大小判断电子态的局域化程度。
%
%### 7.5 PROCAR_OPT
%
%当使用**KPOINTS_OPT**功能时，VASP还会输出**PROCAR_OPT**文件。该文件的格式与PROCAR相同，但包含的是在KPOINTS_OPT中指定的k点处的投影信息。
%
%
%## 九、LOCPOT：局域势文件
%
%### 9.1 LOCPOT概述
%

当需要计算材料的功函数\footnote{因为周期性边界的存在，不同体系的真空能级位置并不相同。通过沿体系表面法线方向对LOCPOT做平面平均，可以得到真空能级位置，这是计算材料功函数的理论依据。}时，就需要\textrm{LOCPOT}文件的输出数据(设置参数\texttt{LVTOT}=\textrm{.TRUE.})。此外，在异质结计算中，通过局域势，可以对齐不同材料的能带。\textrm{LOCPOT}文件保存的是体系的总局域势\textrm{(total local potential)}，\textrm{LOCPOT}文件的内容格式与\textrm{CHGCAR}类似。
%**LOCPOT**文件包含了体系的**总局域势**（total local potential）。其单位为eV。LOCPOT的格式与CHGCAR类似，但**数据排列方式不同**。
%
%### 9.2 LOCPOT的生成条件
%
%要输出LOCPOT文件，必须在INCAR中设置：
%- **LVTOT = .TRUE.** 或 **LVHAR = .TRUE.**
%
%两者的区别在于：
%- `LVTOT = .TRUE.`：输出总局域势 \( V_{\text{LOCPOT}}(\mathbf{r}) = V(\mathbf{r}) + \int \frac{n(\mathbf{r'})}{|\mathbf{r}-\mathbf{r'}|}d\mathbf{r'} + V_{\text{XC}}(\mathbf{r}) \)
%- `LVHAR = .TRUE.`：输出不含交换关联势的Hartree势
%
%### 9.3 LOCPOT的数据组织
%
%LOCPOT的数据组织方式取决于自旋设置：
%- **非自旋极化**（`ISPIN=1`）：单套数据
%- **共线自旋**（`ISPIN=2`）：两套数据（自旋向上和自旋向下）
%- **非共线计算**（`LNONCOLLINEAR = .TRUE.`）：四套数据（旋量表示）
%
%### 9.4 LOCPOT的用途
%
%1. **功函数计算**：通过沿表面法线方向对LOCPOT做平面平均，可以得到真空能级位置，从而计算功函数。
%2. **静电势分析**：分析体系内部的静电势分布。
%3. **能带对齐**：在异质结计算中，通过局域势对齐不同材料的能带。
%
%
%## 十、ELFCAR：电子局域函数文件
如果需要分析``原子之间成键''或考虑``电子局域化程度''等与周期体系中电子局域分布有关的问题，就需要\textrm{ELFCAR}文件中的数据(设置参数\textt{LELF}=\textrm{.TRUE.})，\textrm{ELFCAR}的文件格式与\textrm{CHG}文件类似，因此其存储的电子局域函数\textrm{Electron Localization Function, ELF}可直接用于可视化作图。计算\textrm{ELF}须注意
\begin{itemize}
	\item 为避免圆整误差，建议参数\texttt{PREC}=\textrm{High}；
	\item \textrm{vasp\_nlc}版本尚未实现\textrm{ELF}的计算。
\end{itemize}
%### 10.1 ELFCAR概述
%
%**ELFCAR**文件包含了体系的**电子局域函数**（Electron Localization Function, ELF）。ELF是分析**原子之间成键**和**电子局域化程度**的重要工具。
%
%### 10.2 ELFCAR的生成条件
%
%要输出ELFCAR文件，必须在INCAR中设置：
%```
%LELF = .TRUE.
%```
%
%### 10.3 ELFCAR的格式
%
%ELFCAR采用与**CHGCAR相同的文件格式**：
%- 晶格向量
%- 原子坐标
%- 笛卡尔采样点数 \( N_x, N_y, N_z \)
%- ELF数据值
%
%对于自旋极化计算（`ISPIN=2`），先写入 \( ELF_{\uparrow} \)，再写入 \( ELF_{\downarrow} \)。
%
%### 10.4 ELFCAR的使用注意事项
%
%1. **避免卷绕误差**：建议在INCAR中设置`PREC = High`以避免评估ELFCAR时的卷绕误差。
%2. **非共线计算**：电子局域函数**未实现**于非共线计算。
%3. **可视化**：ELFCAR可通过VESTA等软件进行三维可视化。
%
%
%## 十一、IBZKPT：不可约布里渊区k点文件
%
%### 11.1 IBZKPT概述
%
如果\textrm{KPOINTS}中设置$\vec k$空间网格自动生成，则文件\textrm{IBZKPT}用于保存\textrm{VASP}程序自动产生的不可约布里渊区\textrm{Irreducible Brillouin Zone, IBZ}的$\vec k$点坐标及其权重(由体系点群对称性确定)。
{\fontsize{8.2pt}{6.2pt}\selectfont{
	\lstinputlisting[linerange=1-23]{Figures/IBZKPT} %为保险:~选用文件名绝对路径
}}
\textrm{IBZKPT}文件的内容主要包括:~(1)总的不可约$\vec k$点数目；(2)每个$\vec k$点的坐标及其权重。针对四面体布点积分方案，会给出不等价四面体的数目、四面体的体积和四面体顶点关联列表。
%- 后续行：每个k点的坐标和权重
%- 如果选择了四面体方法，还会包含额外的四面体连接表
%**IBZKPT**文件是在KPOINTS中选择**自动k点网格生成**时由VASP自动产生的输出文件。它包含了**不可约布里渊区**（Irreducible Brillouin Zone）的k点坐标和权重。
%
%### 11.2 IBZKPT的格式
%
%IBZKPT采用与KPOINTS中“手动输入所有k点”相同的格式：
%- 第一行：k点总数
%- 后续行：每个k点的坐标和权重
%- 如果选择了四面体方法，还会包含额外的四面体连接表
%
%### 11.3 IBZKPT的用途
%
%1. **复用k点**：如果多次使用同一套k点网格，可以将IBZKPT复制为KPOINTS以节省时间。
%2. **查看实际k点**：了解VASP实际使用的不可约k点及其权重。
%
%
%## 十二、XDATCAR：离子轨迹文件
%
在结构弛豫和\textrm{MD}计算中，每个离子步的原子径迹(\textrm{trajectory})由\textrmP{XDATCAR}文件记录。参数\texttt{NBLOCK}=\textrm{N}控制每经过\textrm{N}个离子步，在\textrm{XDATCAR}中记录一次当前的原子构型。
{\fontsize{8.2pt}{6.2pt}\selectfont{
	\lstinputlisting[linerange=1-23]{Figures/XDATCAR} %为保险:~选用文件名绝对路径
}}
这些数据可用于可视化结构弛豫或\textrm{MD}过程中原子行经的轨迹，也可以用于提取特定离子步的原子构型，用于后续分析。
%### 12.1 XDATCAR概述
%
%**XDATCAR**文件包含了**每个离子步的原子位置**。它是追踪结构弛豫或分子动力学过程中原子运动轨迹的核心文件。
%
%### 12.2 XDATCAR的生成
%
%XDATCAR由`NBLOCK`标签控制写入频率——每经过`NBLOCK`个离子步，当前的离子构型就会被写入XDATCAR。
%
%### 12.3 XDATCAR的用途
%
%1. **轨迹分析**：可视化原子在弛豫或MD过程中的运动轨迹。
%2. **结构快照**：提取特定离子步的结构用于后续分析。
%3. **OVITO可视化**：XDATCAR可直接导入OVITO等可视化软件进行动画展示。
%
%
\textrm{DFT}计算的核心数据无非是体系总能量(总自由能)、电子结构(含态密度和能带)、电荷密度与\textrm{Kohn-Sham}本征态波函数。对应到\textrm{VASP}计算就是\textrm{OSZICAR}、\textrm{DOSCAR}、\textrm{EIGENVAL}、\textrm{CHGCAR}和\textrm{WAVECAR}。此外，上面还介绍了一些常用于数据分析和可视化的输出文件。还有一些更少用到的文件，这里简要提一下
\begin{itemize}
	\item \textrm{PARCHG}文件存储的是分数电荷密度\textrm{(partial charge densities)}，通常用于可视化特定能量范围或特定$\vec k$点、能带附近的电荷密度分布。
	\item \textrm{PCDAT}文件存储的是\textrm{MD}计算中的对相关函数\textrm{(pair correlation function)}，主要用于\textrm{MD}模拟中的径向分布函数分析。
	\item \textrm{REPORT}文件记录的是各种\textrm{MD}计算的输出结果，如伞状积分\textrm{(umbrella integration)}的数值等。
\end{itemize}
对于初学\textrm{VASP}的用户，建议尽早熟悉\textrm{OUTCAR}文件，并能逐渐旁及\textrm{OSZICAR}、\textrm{DOSCAR}、\textrm{EIGENVAL}、\textrm{CHGCAR}等文件。
%%
%## 十三、其他输出文件
%
%### 13.1 CHG文件
%
%**CHG**文件包含电荷密度、晶格向量和原子坐标。它与CHGCAR的区别在于**不包含PAW一中心占据数**，主要用于**可视化**。
%
%### 13.2 PARCHG文件
%
%**PARCHG**文件包含**部分电荷密度**（partial charge densities）。它通常用于可视化特定能量范围或特定k点、能带的电荷密度分布。
%
%### 13.3 PCDAT文件
%
%**PCDAT**文件包含**对关联函数**（pair correlation function）。它主要用于分子动力学模拟中的径向分布函数分析。
%
%### 13.4 REPORT文件
%
%**REPORT**文件包含了各种分子动力学计算的输出结果，如伞形积分（umbrella integration）等。
%
%### 13.5 BSEFATBAND文件
%
%**BSEFATBAND**文件包含BSE（Bethe-Salpeter Equation）的本征值，用于绘制“fatband”图。
%
%### 13.6 TMPCAR文件
%
%**TMPCAR**文件包含上一步离子步的波函数和离子位置。
%
%### 13.7 PROCAR_OPT
%
%当使用**KPOINTS_OPT**功能时，VASP会输出**PROCAR_OPT**文件，其格式与PROCAR相同。
%
%
%%## 十四、现代输出格式：vasprun.xml与vaspout.h5
%%
%%### 14.1 vasprun.xml
%%
%%**vasprun.xml**是VASP的**XML格式**主输出文件。它包含了与OUTCAR类似的信息，但以结构化的XML格式呈现。
%%
%%vasprun.xml的优势在于：
%%- **易于解析**：XML格式便于程序化读取
%%- **与p4vasp兼容**：p4vasp可以直接读取vasprun.xml进行可视化
%%
%%### 14.2 vaspout.h5
%%
%%**vaspout.h5**是VASP的**HDF5格式**主输出文件。当启用HDF5支持时，VASP会生成此文件。
%%
%%vaspout.h5的优势在于：
%%- **高效存储**：HDF5格式支持大规模数据的高效存储和读取
%%- **与py4vasp兼容**：py4vasp后处理工具需要vaspout.h5作为输入
%%
%%### 14.3 vaspwave.h5
%%
%%当输出重定向到HDF5时，**vaspwave.h5**文件包含电荷密度和波函数。
%%
%
%## 十五、输出文件的协同使用
%
%VASP的各类输出文件并非孤立存在，它们之间存在着紧密的协同关系：
%
%**1. 结构信息的流转**：POSCAR → CONTCAR → 下一轮计算的POSCAR
%
%**2. 波函数与电荷密度的续算**：WAVECAR和CHGCAR可用于加速后续计算的收敛
%
%**3. MD中的一致性**：在分子动力学中，WAVECAR、CHGCAR和CONTCAR是相互兼容的，可以一致地用于续算
%
%**4. 后处理的输入链**：DOSCAR + OUTCAR + POSCAR → 态密度图绘制
%
%**5. 能带分析的组合**：EIGENVAL + KPOINTS + OUTCAR → 能带结构图
%
%
%## 十六、输出文件的实践指南
%
%### 16.1 计算完成后的检查清单
%
%1. **检查OUTCAR**：确认`reached required accuracy`或类似信息，确保电子步和离子步均已收敛。
%2. **检查OSZICAR**：快速确认最终能量和收敛状态。
%3. **检查CONTCAR**：确认结构文件已正确生成。
%4. **检查所需性质文件**：根据需要确认DOSCAR、EIGENVAL、CHGCAR等是否已生成。
%
%### 16.2 存储管理建议
%
%1. **大文件控制**：对于大型体系，WAVECAR可能达到数GB。如不需要，设置`LWAVE = .FALSE.`。
%2. **定期清理**：计算完成后，可将CHGCAR和WAVECAR压缩存档或删除以节省空间。
%3. **保留关键文件**：OUTCAR、CONTCAR、OSZICAR通常需要长期保留。
%
%### 16.3 后处理工具推荐
%
%| 工具 | 用途 |
%|------|------|
%| **p4vasp** | DOS/PDOS绘制、结构可视化 |
%| **VASPKIT** | 提取能量、生成能带图、处理DOSCAR等 |
%| **VESTA** | CHGCAR/ELFCAR三维可视化 |
%| **OVITO** | XDATCAR轨迹动画 |
%| **PyProcar** | PROCAR的Python分析工具包 |
%
%### 16.4 常见问题与解决方案
%
%1. **OUTCAR中未找到能量**：检查计算是否正常结束，或搜索`TOTEN`和`entropy=`两个关键词。
%2. **DOSCAR为空或异常**：检查INCAR中的`NEDOS`、`EMIN`、`EMAX`设置是否合理。
%3. **CONTCAR与POSCAR完全相同**：对于静态计算（`NSW=0`）这是正常的；对于结构优化，检查`IBRION`和`NSW`设置是否正确。
%4. **PROCAR未生成**：检查`LORBIT`设置和`NPAR`值（并行时需`NPAR=1`）。
%
%
%## 十七、总结
%
%VASP的输出文件体系是一个完整而精密的信息系统，每个文件都承载着特定类型的计算结果。**OUTCAR**作为最核心的输出文件，提供了从输入参数到计算结果的全部详细日志；**CONTCAR**保存了更新后的结构，是后续计算的起点；**DOSCAR**和**EIGENVAL**分别提供了态密度和能带信息；**CHGCAR**和**WAVECAR**存储了电荷密度和波函数；**PROCAR**则揭示了电子态的轨道成分。
%
%理解每个输出文件的格式、内容和用途，是准确解读VASP计算结果的基础。初学者应从最常用的OUTCAR、CONTCAR、DOSCAR和EIGENVAL开始，逐步扩展到其他专业文件。在实践中，建议养成以下习惯：
%
%1. **每次计算后检查OUTCAR**，确认收敛状态。
%2. **保存CONTCAR**作为最终结构，并复制为POSCAR用于后续计算。
%3. **根据研究目的选择性地保留**CHGCAR、WAVECAR等大文件。
%4. **善用后处理工具**，将原始数据转化为可视化图表和可发表的结果。
%
%正如VASP官方所建议的：始终保留一份OSZICAR文件的副本，它可能包含重要的信息。同样，对每一个输出文件都应保持同样的重视——因为最终的科学发现，往往就隐藏在这些看似平凡的ASCII字符和二进制数据之中。
